\chapter{Requisitos del sistema}
\label{annex:rf}

\section{Requisitos Funcionales}

\subsection{RF-01: Generación del Árbol de Sintaxis Abstracta (AST)}

\textbf{Descripción:}
El sistema debe procesar el programa ingresado por el usuario utilizando la gramática previamente definida para generar su correspondiente Árbol de Sintaxis Abstracta (AST).

\textbf{Entradas:}
\begin{itemize}
    \item Gramática BNF definida y confirmada por el usuario en el modal de gramática.
    \item Programa ingresado como expresión de prueba en la consola del editor.
\end{itemize}

\textbf{Prioridad:} Alta

\textbf{Criterios de aceptación:}
\begin{itemize}
    \item El sistema invoca el intérprete Racket con los archivos del editor y la expresión de entrada; el AST resultante es serializado en formato JSON y transmitido al cliente.
    \item El AST generado es renderizado en el panel ASTViewer reflejando la estructura jerárquica del programa evaluado.
    \item Cuando el programa contiene errores sintácticos o de evaluación, el sistema muestra un mensaje de error en la consola con el prefijo \texttt{✕} y nivel visual diferenciado en rojo; el panel ASTViewer conserva el último AST válido.
    \item El AST corresponde fielmente a la última expresión evaluada exitosamente.
\end{itemize}

\subsection{RF-02: Visualización gráfica e interactiva del AST}

\textbf{Descripción:}
El sistema debe mostrar el árbol de sintaxis abstracta generado de forma gráfica e interactiva, permitiendo la exploración visual de su estructura.

\textbf{Entradas:}
\begin{itemize}
    \item AST serializado en formato JSON producido por el intérprete Racket.
\end{itemize}

\textbf{Prioridad:} Alta

\textbf{Criterios de aceptación:}
\begin{itemize}
    \item El panel ASTViewer renderiza la representación jerárquica del AST con nodos que indican el tipo de dato retornado por el intérprete.
    \item La visualización se actualiza automáticamente tras cada ejecución exitosa, sin necesidad de acción adicional por parte del usuario.
    \item El panel es responsivo y se adapta al tamaño del contenedor asignado dentro del layout redimensionable.
    \item En modo paso a paso, el panel refleja el AST correspondiente al paso de evaluación actualmente visualizado.
    \item No se presentan errores de renderizado ni contenido huérfano durante el ciclo de vida de la sesión.
\end{itemize}

\subsection{RF-03: Ejecución y evaluación del programa}

\textbf{Descripción:}
El sistema debe permitir ejecutar el programa del usuario e interpretar su resultado, siguiendo el proceso de evaluación definido por la gramática y el intérprete generado.

\textbf{Entradas:}
\begin{itemize}
    \item Archivos del editor (\texttt{grammar.rkt}, \texttt{environment.rkt}, \texttt{main.rkt}, \texttt{utils.rkt}).
    \item Expresión de prueba ingresada por el usuario en la consola.
\end{itemize}

\textbf{Prioridad:} Alta

\textbf{Criterios de aceptación:}
\begin{itemize}
    \item El sistema transmite los archivos del editor y la expresión de entrada a la ruta \texttt{POST /api/run} y recibe el arreglo de pasos de evaluación.
    \item El resultado de cada paso de evaluación es presentado en la consola con el nivel de salida correspondiente: verde (\texttt{output}) para valores, rojo (\texttt{error}) para excepciones capturadas, gris (\texttt{info}) para resultados \texttt{void}.
    \item Los errores producidos por \texttt{eopl:error} son capturados por el runtime, limpiados de trazas internas de Racket, y presentados con el prefijo \texttt{✕} en color rojo.
    \item La ejecución puede ser cancelada por el usuario durante el proceso; el servidor aborta el subproceso Racket correspondiente.
    \item El tiempo máximo de ejecución del subproceso Racket es de 15 segundos; transcurrido ese tiempo, el sistema reporta el error de tiempo límite en la consola.
    \item El sistema soporta un modo de evaluación paso a paso que permite avanzar la ejecución de forma controlada mostrando el estado del AST y del ambiente en cada paso.
\end{itemize}

\subsection{RF-04: Representación del ambiente de ejecución}

\textbf{Descripción:}
El sistema debe representar el ambiente de evaluación como una estructura jerárquica de frames que se actualiza en cada ejecución, permitiendo al estudiante observar la evolución del estado del programa.

\textbf{Entradas:}
\begin{itemize}
    \item Instantáneas del ambiente de evaluación capturadas por el módulo \texttt{\_tracking.rkt} durante la ejecución de cada paso.
\end{itemize}

\textbf{Prioridad:} Alta

\textbf{Criterios de aceptación:}
\begin{itemize}
    \item El panel EnvironmentPanel muestra el ambiente de evaluación como una pila de frames jerárquicos, donde cada frame lista los bindings activos con su nombre, valor serializado y tipo de dato inferido.
    \item El ambiente mostrado corresponde al último paso de evaluación completado; en modo paso a paso, se actualiza con cada avance del usuario.
    \item La estructura de frames refleja fielmente la cadena de entornos construida durante la evaluación, permitiendo al estudiante identificar el alcance léxico de cada variable.
    \item El panel permanece visible y no colapsa ante la ausencia de bindings en el ambiente actual.
\end{itemize}

\subsection{RF-05: Biblioteca de ejemplos}

\textbf{Descripción:}
El sistema debe ofrecer una biblioteca de ejemplos predefinidos orientados a los indicadores de logro identificados como base conceptual del prototipo, con el propósito de facilitar el aprendizaje mediante casos de uso concretos.

\textbf{Entradas:}
\begin{itemize}
    \item Selección de ejemplo por parte del usuario en el modal de bienvenida.
\end{itemize}

\textbf{Prioridad:} Media

\textbf{Criterios de aceptación:}
\begin{itemize}
    \item Los ejemplos son cargados desde el sistema de archivos del servidor durante el renderizado inicial (Server Component) y transmitidos al cliente como propiedades del componente orquestador.
    \item El sistema presenta la lista de ejemplos disponibles en el modal de bienvenida con nombre y descripción breve de cada uno.
    \item Al seleccionar un ejemplo, los archivos del editor son reemplazados por los archivos correspondientes al ejemplo seleccionado sin pérdida de estado de la interfaz.
    \item Los ejemplos cubren los cinco indicadores de logro seleccionados en el diagnóstico: IL 1.1.6, IL 1.2.3, IL 1.2.6, IL 1.2.7 e IL 1.3.1.
    \item Cada ejemplo es compatible con el pipeline de gramática y puede ser ejecutado sin modificaciones adicionales.
\end{itemize}

\subsection{RF-06: Soporte y validación de sintaxis en Racket}

\textbf{Descripción:}
El sistema debe ser capaz de reconocer la sintaxis del lenguaje Racket en el editor y validar la gramática BNF ingresada por el usuario, reportando errores de forma clara y oportuna.

\textbf{Entradas:}
\begin{itemize}
    \item Texto BNF ingresado en el modal de gramática.
    \item Código fuente Racket presente en los archivos del editor.
\end{itemize}

\textbf{Prioridad:} Alta

\textbf{Criterios de aceptación:}
\begin{itemize}
    \item El editor Monaco provee resaltado de sintaxis para archivos con extensión \texttt{.rkt}.
    \item El pipeline de gramática detecta y reporta errores léxicos en la especificación BNF durante el procesamiento; los mensajes de error son presentados en el modal de gramática antes de generar los archivos.
    \item El sistema no genera los archivos del intérprete ni habilita la ejecución cuando la gramática BNF contiene errores léxicos o sintácticos.
    \item Los errores de evaluación producidos por el intérprete Racket son presentados en la consola con nivel de error y sin exponer trazas internas del runtime de Racket al usuario.
\end{itemize}

\subsection{RF-07: Visualización de los componentes internos del intérprete}

\textbf{Descripción:}
El sistema debe permitir al estudiante visualizar los archivos que componen el intérprete generado, mostrando su estructura de forma organizada con distinción entre las secciones técnicas y las secciones destinadas a la implementación por parte del estudiante.

\textbf{Entradas:}
\begin{itemize}
    \item Archivos generados por el pipeline de gramática (\texttt{grammar.rkt}, \texttt{environment.rkt}, \texttt{main.rkt}, \texttt{utils.rkt}).
\end{itemize}

\textbf{Prioridad:} Media

\textbf{Criterios de aceptación:}
\begin{itemize}
    \item El panel de edición organiza los archivos del intérprete en pestañas independientes, accesibles mediante navegación por clic.
    \item Las líneas generadas automáticamente por el pipeline (código de infraestructura del intérprete) son marcadas como de solo lectura en el editor y presentadas con distinción visual respecto a las líneas editables.
    \item Las líneas destinadas a la implementación por parte del estudiante (stubs de \texttt{eval-program}, \texttt{eval-expression}, \texttt{apply-primitive}) permanecen activas y editables.
    \item No se permite la modificación de las líneas protegidas durante la sesión de trabajo.
\end{itemize}

\subsection{RF-08: Plantilla base editable y exportación del entorno}

\textbf{Descripción:}
El sistema debe proporcionar una plantilla base del intérprete donde el estudiante pueda completar la implementación, y permitir la descarga del proyecto completo en formato \texttt{.zip} compatible con la librería EOPL.

\textbf{Entradas:}
\begin{itemize}
    \item Gramática BNF confirmada por el usuario.
    \item Código editado por el usuario en las secciones habilitadas del editor.
\end{itemize}

\textbf{Prioridad:} Alta

\textbf{Criterios de aceptación:}
\begin{itemize}
    \item Al confirmar la gramática, el pipeline genera automáticamente los cuatro archivos del intérprete con los stubs de las funciones a implementar, cargándolos en el editor sin intervención adicional del usuario.
    \item Las secciones de infraestructura técnica del intérprete están protegidas y no son modificables por el usuario.
    \item Las secciones de implementación están identificadas con comentarios \texttt{;; TODO} que orientan al estudiante sobre qué debe completar.
    \item El sistema genera un archivo \texttt{.zip} que contiene todos los archivos del editor en su estado actual, incluyendo las modificaciones realizadas por el usuario.
    \item El archivo descargado incluye una estructura de directorios compatible con la librería EOPL y puede ser abierto directamente en DrRacket.
\end{itemize}

\section{Requisitos No Funcionales}

\subsection{RNF-01: Usabilidad}

\textbf{Descripción:}
El sistema debe ser intuitivo y operable por estudiantes sin experiencia avanzada en compiladores, minimizando la carga cognitiva extrínseca impuesta por la interfaz.

\textbf{Prioridad:} Alta

\textbf{Criterios de aceptación:}
\begin{itemize}
    \item Los usuarios pueden iniciar una sesión de trabajo, definir una gramática y ejecutar un programa sin consultar documentación externa.
    \item Los mensajes de error presentados en la consola son comprensibles sin conocimiento del funcionamiento interno de Racket.
    \item La distribución de paneles sigue un orden visual coherente con el flujo de trabajo del estudiante: definición de gramática, edición del intérprete, ejecución y observación de resultados.
    \item Las secciones del editor que no deben ser modificadas están diferenciadas visualmente de forma que el estudiante pueda identificarlas sin instrucción explícita.
\end{itemize}

\subsection{RNF-02: Rendimiento}

\textbf{Descripción:}
El sistema debe completar la generación del AST y la evaluación del programa en un tiempo que no interrumpa el flujo cognitivo del estudiante durante la exploración.

\textbf{Prioridad:} Alta

\textbf{Criterios de aceptación:}
\begin{itemize}
    \item La transformación BNF a archivos Racket (pipeline de gramática) se completa en el cliente en menos de 1 segundo para gramáticas de hasta 20 producciones.
    \item La respuesta de la ruta \texttt{/api/run} para programas de complejidad habitual en la asignatura FLP se recibe en menos de 5 segundos bajo condiciones normales de red.
    \item El tiempo máximo de ejecución del subproceso Racket está limitado a 15 segundos; transcurrido ese tiempo, el sistema notifica al usuario y libera los recursos asociados.
    \item No se producen bloqueos del hilo principal del navegador durante la ejecución de la petición al servidor.
\end{itemize}

\subsection{RNF-03: Accesibilidad web}

\textbf{Descripción:}
El sistema debe ejecutarse en navegadores web modernos sin requerir instalación de software adicional, plugins ni autenticación por parte del usuario.

\textbf{Prioridad:} Alta

\textbf{Criterios de aceptación:}
\begin{itemize}
    \item La aplicación carga y opera correctamente en las versiones actuales de Google Chrome, Mozilla Firefox y Microsoft Edge.
    \item No se requiere instalación de extensiones, plugins ni entornos de ejecución adicionales en el equipo del usuario.
    \item El acceso al sistema no requiere registro ni autenticación de ningún tipo.
    \item La interfaz es funcional en resoluciones de pantalla iguales o superiores a 1280 × 720 píxeles.
\end{itemize}

\subsection{RNF-04: Mantenibilidad}

\textbf{Descripción:}
El código fuente debe estar organizado de forma modular, con separación explícita de responsabilidades, para facilitar la incorporación de mejoras futuras.

\textbf{Prioridad:} Media

\textbf{Criterios de aceptación:}
\begin{itemize}
    \item Las definiciones de tipos e interfaces TypeScript residen exclusivamente en el directorio \texttt{app/types/}, con importaciones explícitas desde los módulos que las consumen.
    \item Las funciones puras y utilitarios que no dependen del estado React residen en \texttt{app/lib/} y no en los archivos de componentes.
    \item El pipeline de gramática está segmentado en módulos independientes (lexer, parser, generadores) con responsabilidades no solapadas.
    \item Es posible añadir un nuevo generador de archivo Racket incorporando un módulo en \texttt{app/lib/} sin modificar los componentes de interfaz.
\end{itemize}

\subsection{RNF-05: Escalabilidad}

\textbf{Descripción:}
La arquitectura del sistema debe permitir incorporar nuevos ejemplos, gramáticas de mayor complejidad o características adicionales sin afectar la estabilidad de las funcionalidades existentes.

\textbf{Prioridad:} Media

\textbf{Criterios de aceptación:}
\begin{itemize}
    \item La incorporación de nuevos ejemplos se realiza añadiendo archivos al directorio de ejemplos del servidor, sin modificar código de componentes ni de la API.
    \item El pipeline de gramática procesa gramáticas con hasta 20 producciones sin degradación observable del rendimiento en el cliente.
    \item La adición de un nuevo panel de visualización puede realizarse como un componente React independiente integrable en el layout sin alterar el comportamiento de los paneles existentes.
    \item El sistema de contenedores Docker permite escalar el servicio de ejecución de forma independiente al servidor de la interfaz.
\end{itemize}
